\documentclass[../../main.tex]{subfiles}% 注意这里的文件路径不能用 ./main.tex ,否则用latexmk编译子文件会报错
\graphicspath{{\subfix{./image/}}{\subfix{../../image/}}} % 指定图片引用路径,后续可以直接使用图片文件名
% 如果图片存放在上级目录,则使用 ../../image/ 作为备用路径

% 例如:
% \begin{figure}[H]
% \centering
% \includegraphics[width=0.3\textwidth]{图.png}
% \caption{}
% \label{figure:图}
% \end{figure}
% 注意:上述\label{}一定要放在\caption{}之后,否则引用图片序号会只会显示??.

\begin{document}

% SOURCE p=25; printed=10; object=ch00-sec05-title; region=lower
\section{扩展实数系}\label{section:ch00-extended-real-numbers}

% SOURCE p=25; printed=10; object=ch00-extended-real-definition; region=bottom

\begin{definition}[扩展实数系及其确界]\label{definition:ch00-extended-real-definition}
常常有必要向 $\mathbb{R}$ 添加两个额外的点 $\infty$（即 $+\infty$）和 $-\infty$，构成\textbf{扩展实数系} $\overline{\mathbb{R}}=\mathbb{R}\cup\{-\infty,\infty\}$，并规定对所有 $x\in\mathbb{R}$ 都有 $-\infty<x<\infty$，从而扩展 $\mathbb{R}$ 上通常的序。于是，$\mathbb{R}$ 的完备性可表述为：$\overline{\mathbb{R}}$ 的每个子集 $A$ 都有最小上界，即\textbf{上确界}，以及最大下界，即\textbf{下确界}，分别记为 $\sup A$ 和 $\inf A$。若 $A=\{a_1,\ldots,a_n\}$，也记
\begin{align*}
\max(a_1,\ldots,a_n)=\sup A,\qquad \min(a_1,\ldots,a_n)=\inf A.
\end{align*}
\end{definition}


% SOURCE p=26; printed=11; object=ch00-limsup-liminf; region=top

\begin{definition}[序列的上极限与下极限]\label{definition:ch00-limsup-liminf}
由完备性可知，$\overline{\mathbb{R}}$ 中每个序列 $\{x_n\}$ 都有\textbf{上极限}和\textbf{下极限}：
\begin{align*}
\limsup x_n=\inf_{k\geqslant1}\Bigl(\sup_{n\geqslant k}x_n\Bigr),\qquad
\liminf x_n=\sup_{k\geqslant1}\Bigl(\inf_{n\geqslant k}x_n\Bigr).
\end{align*}
\end{definition}


\begin{remark}
序列 $\{x_n\}$（在 $\mathbb{R}$ 中）收敛，当且仅当这两个数相等（且有限），此时其极限就是这个共同值。
\end{remark}

\begin{definition}[函数的上极限]\label{definition:ch00-function-limsup}
也可以对函数 $f:\mathbb{R}\to\overline{\mathbb{R}}$ 定义 $\limsup$ 和 $\liminf$，例如：
\begin{align*}
\limsup_{x\to a}f(x)=\inf_{\delta>0}\Bigl(\sup_{0<|x-a|<\delta}f(x)\Bigr).
\end{align*}
\end{definition}


% SOURCE p=26; printed=11; object=ch00-extended-arithmetic; region=upper-middle

\begin{definition}[扩展实数的运算约定]\label{definition:ch00-extended-arithmetic}
$\mathbb{R}$ 上的算术运算可以部分地扩展到 $\overline{\mathbb{R}}$：
\begin{align*}
x\pm\infty&=\pm\infty\quad(x\in\mathbb{R}),&
\infty+\infty&=\infty,&-\infty-\infty&=-\infty,\\
x\cdot(\pm\infty)&=\pm\infty\quad(x>0),&
x\cdot(\pm\infty)&=\mp\infty\quad(x<0).
\end{align*}
我们不定义 $\infty-\infty$，但约定，除非另有说明，
\begin{align*}
0\cdot(\pm\infty)=0.
\end{align*}
\end{definition}


\begin{remark}
（在测度论中，表达式 $0\cdot\infty$ 不时出现；出于各种原因，它几乎总应被解释为 $0$。）
\end{remark}


% SOURCE p=26; printed=11; object=ch00-interval-notation; region=middle

\begin{definition}[扩展实数中的区间]\label{definition:ch00-interval-notation}
对于 $\overline{\mathbb{R}}$ 中的区间，采用下列记号：若 $-\infty\leqslant a<b\leqslant\infty$，则
\begin{align*}
(a,b)&=\{x:a<x<b\},& [a,b]&=\{x:a\leqslant x\leqslant b\},\\
(a,b]&=\{x:a<x\leqslant b\},& [a,b)&=\{x:a\leqslant x<b\}.
\end{align*}
\end{definition}


% SOURCE p=26; printed=11; object=ch00-uncountable-sum-definition; region=lower-middle
我们偶尔会遇到非负数的不可数和。
\begin{definition}[非负数的任意和]\label{definition:ch00-uncountable-sum-definition}
若 $X$ 为任意集合，$f:X\to[0,\infty]$，则将 $\sum_{x\in X}f(x)$ 定义为其有限部分和的上确界：
\begin{align*}
\sum_{x\in X}f(x)=\sup\Bigl\{\sum_{x\in F}f(x):F\subseteq X,\ F\text{ 为有限集}\Bigr\}.
\end{align*}
\end{definition}

（以后我们会认识到，这就是 $f$ 关于 $X$ 上计数测度的积分。）

% SOURCE p=26; printed=11; object=ch00-nonnegative-sums; number=0.20; region=lower
\begin{proposition}\label{proposition:ch00-nonnegative-sums}
给定 $f:X\to[0,\infty]$，令 $A=\{x:f(x)>0\}$。若 $A$ 不可数，则 $\sum_{x\in X}f(x)=\infty$。若 $A$ 为可数无限集，则 $\sum_{x\in X}f(x)=\sum_{1}^{\infty}f(g(n))$，其中 $g:\mathbb{N}\to A$ 是任意双射，右端的和是通常的无穷级数。
\end{proposition}

% SOURCE p=26; printed=11; object=ch00-proof-nonnegative-sums-start; region=bottom
\begin{proof}
有 $A=\bigcup_{1}^{\infty}A_n$，其中 $A_n=\{x:f(x)>1/n\}$。若 $A$ 不可数，则必有某个 $A_n$ 不可数；且对于 $A_n$ 的有限子集 $F$，有 $\sum_{x\in F}f(x)>\card(F)/n$，由此得到 $\sum_{x\in X}f(x)=\infty$。
% SOURCE p=27; printed=12; object=ch00-proof-nonnegative-sums-end; region=top
若 $A$ 为可数无限集，$g:\mathbb{N}\to A$ 是双射，且 $B_N=g(\{1,\ldots,N\})$，则 $A$ 的每个有限子集 $F$ 都包含在某个 $B_N$ 中。因此
\begin{align*}
\sum_{x\in F}f(x)\leqslant\sum_{1}^{N}f(g(n))\leqslant\sum_{x\in X}f(x).
\end{align*}
对 $N$ 取上确界，得到
\begin{align*}
\sum_{x\in F}f(x)\leqslant\sum_{1}^{\infty}f(g(n))\leqslant\sum_{x\in X}f(x),
\end{align*}
再对 $F$ 取上确界，即得所需结论。
\end{proof}

% SOURCE p=27; printed=12; object=ch00-monotone-functions; region=upper-middle
下面给出关于（扩展）实值函数的一些术语。数之间的关系应用到函数时，理解为逐点成立。因此，$f\leqslant g$ 表示对每个 $x$ 都有 $f(x)\leqslant g(x)$，而 $\max(f,g)$ 是在 $x$ 处取值为 $\max(f(x),g(x))$ 的函数。
\begin{definition}[单调函数]\label{definition:ch00-monotone-functions}
若 $X\subseteq\overline{\mathbb{R}}$ 且 $f:X\to\overline{\mathbb{R}}$，当 $x\leqslant y$ 总能推出 $f(x)\leqslant f(y)$ 时，称 $f$ \textbf{递增}；当 $x<y$ 总能推出 $f(x)<f(y)$ 时，称 $f$ \textbf{严格递增}；\textbf{递减}的定义与此类似。递增或递减的函数称为\textbf{单调函数}。
\end{definition}


% SOURCE p=27; printed=12; object=ch00-one-sided-limits; region=middle

\begin{proposition}\label{proposition:ch00-one-sided-limits}
若 $f:\mathbb{R}\to\mathbb{R}$ 为递增函数，则 $f$ 在每一点都有右极限和左极限：
\begin{align*}
f(a+)&=\lim_{x\searrow a}f(x)=\inf_{x>a}f(x),&
f(a-)&=\lim_{x\nearrow a}f(x)=\sup_{x<a}f(x).
\end{align*}
此外，极限值 $f(\infty)=\sup_{x\in\mathbb{R}}f(x)$ 和 $f(-\infty)=\inf_{x\in\mathbb{R}}f(x)$ 存在（可能等于 $\pm\infty$）。
\end{proposition}

\begin{definition}[右连续与左连续]\label{definition:ch00-one-sided-continuity}
若对所有 $a\in\mathbb{R}$ 都有 $f(a)=f(a+)$，则称 $f$ \textbf{右连续}；若对所有 $a\in\mathbb{R}$ 都有 $f(a)=f(a-)$，则称 $f$ \textbf{左连续}。
\end{definition}


% SOURCE p=27; printed=12; object=ch00-absolute-value-norm; region=middle-lower

\begin{definition}[绝对值、模与 Euclid 范数]\label{definition:ch00-absolute-value-norm}
对于 $\mathbb{R}$ 或 $\mathbb{C}$ 中的点 $x$，$|x|$ 表示 $x$ 通常的绝对值或模，$|a+\mathrm{i}b|=\sqrt{a^2+b^2}$。对于 $\mathbb{R}^n$ 或 $\mathbb{C}^n$ 中的点 $x$，$|x|$ 表示 Euclid 范数：
\begin{align*}
|x|=\Bigl[\sum_{1}^{n}|x_j|^2\Bigr]^{1/2}.
\end{align*}
\end{definition}


% SOURCE p=27; printed=12; object=ch00-open-real-set; region=lower

\begin{definition}[实数集中的开集]\label{definition:ch00-open-real-set}
回忆，集合 $U\subseteq\mathbb{R}$ 为\textbf{开集}，是指对每个 $x\in U$，$U$ 都包含一个以 $x$ 为中心的区间。
\end{definition}


% SOURCE p=27; printed=12; object=ch00-open-interval-decomposition; number=0.21; region=lower
\begin{proposition}\label{proposition:ch00-open-interval-decomposition}
$\mathbb{R}$ 中每个开集都是可数个两两不交的开区间之并。
\end{proposition}

% SOURCE p=27; printed=12; object=ch00-proof-open-interval-decomposition; region=bottom
\begin{proof}
若 $U$ 为开集，对每个 $x\in U$，考虑所有满足 $x\in I\subseteq U$ 的开区间 $I$ 所组成的族 $\mathscr{I}_x$。容易检查，任意一个具有公共点的开区间族的并仍是开区间，因此 $J_x=\bigcup_{I\in\mathscr{I}_x}I$ 是开区间；它是 $\mathscr{I}_x$ 的最大元。若 $x,y\in U$，则或者 $J_x=J_y$，或者 $J_x\cap J_y=\varnothing$，否则 $J_x\cup J_y$ 将是 $\mathscr{I}_x$ 中比 $J_x$ 更大的开区间。因此，若 $\mathscr{J}=\{J_x:x\in U\}$，则 $\mathscr{J}$ 中的不同成员两两不交，且 $U=\bigcup_{J\in\mathscr{J}}J$。对每个 $J\in\mathscr{J}$，选取有理数 $f(J)\in J$。这样定义的映射 $f:\mathscr{J}\to\mathbb{Q}$ 是单射，因为若 $J\ne J'$，则 $J\cap J'=\varnothing$；因此 $\mathscr{J}$ 可数。
\end{proof}
\end{document}
